\documentclass[10pt,a4paper]{article}
\usepackage[a4paper,margin=14mm,top=12mm,bottom=14mm]{geometry}
\usepackage{fontspec}
\setmainfont{Liberation Sans}
\usepackage{polyglossia}
\setdefaultlanguage{french}
\usepackage{amsmath}
\usepackage{xcolor}
\usepackage{graphicx}
\usepackage[most]{tcolorbox}
\usepackage{tabularx,booktabs,array,multirow}
\usepackage{enumitem}
\usepackage{fancyhdr}
\usepackage{lastpage}
\usepackage[hidelinks]{hyperref}
\definecolor{uimm}{HTML}{005677}
\definecolor{cyan}{HTML}{13A2DD}
\definecolor{fond}{HTML}{E6F2FA}
\definecolor{ok}{HTML}{2E7D32}
\definecolor{ko}{HTML}{B3261E}
\definecolor{violet}{HTML}{7B1FA2}
\setlength{\parindent}{0pt}
\setlength{\parskip}{3pt}
\setlist{nosep}
\setlist[itemize]{leftmargin=13pt,label=\textcolor{cyan}{\textbullet}}
\newcommand{\titre}[1]{\vspace{5pt}{\color{uimm}\large\bfseries #1}\par\vspace{-3pt}{\color{cyan}\rule{\linewidth}{1.2pt}}\par\vspace{0pt}}
\tcbset{boite/.style={colback=fond,colframe=uimm,boxrule=0.6pt,arc=2pt,left=6pt,right=6pt,top=4pt,bottom=4pt},
        defi/.style={colback=violet!6,colframe=violet,boxrule=0.8pt,arc=2pt,left=6pt,right=6pt,top=4pt,bottom=4pt,
                     title={\bfseries Défi},coltitle=white,colbacktitle=violet,fonttitle=\small}}
\newcommand{\bandeau}[3]{%
\begin{tcolorbox}[colback=uimm,colframe=uimm,arc=2pt,left=4pt,right=8pt,top=4pt,bottom=4pt]
\begin{minipage}[c]{18mm}\setlength{\fboxsep}{1.2pt}\fcolorbox{white}{white}{\includegraphics[height=16.5mm]{logo_uimm.jpg}}\end{minipage}\hspace{4mm}%
\begin{minipage}[c]{\dimexpr\linewidth-22mm\relax}\color{white}
{\small #1}\\[1pt]{\LARGE\bfseries #2}\\[2pt]{\large #3}\hfill{\small CQPM Technicien qualité}
\end{minipage}
\end{tcolorbox}}
\newcommand{\reponse}[1][1]{\par\vspace{2pt}{\color{gray!60}\foreach \i in {1,...,#1}{\rule{\linewidth}{0.3pt}\par\vspace{9pt}}}}
\usepackage{pgffor}
\pagestyle{fancy}\fancyhf{}
\renewcommand{\headrulewidth}{0pt}
\newcommand{\pied}[1]{\fancyfoot[L]{\footnotesize\color{uimm}#1}\fancyfoot[R]{\footnotesize\color{uimm}Page \thepage{} sur \pageref{LastPage}}}

\usepackage{tikz}
\usepackage{longtable}
\renewcommand{\bandeau}[3]{%
\begin{tcolorbox}[colback=uimm,colframe=uimm,arc=2pt,left=4pt,right=8pt,top=4pt,bottom=4pt]
\begin{minipage}[c]{18mm}\setlength{\fboxsep}{1.2pt}\fcolorbox{white}{white}{\includegraphics[height=16.5mm]{logo_uimm.jpg}}\end{minipage}\hspace{4mm}%
\begin{minipage}[c]{\dimexpr\linewidth-22mm\relax}\color{white}
{\small #1}\\[1pt]{\LARGE\bfseries #2}\\[2pt]{\large #3}
\end{minipage}
\end{tcolorbox}}
% en-tête d'une fiche : code, titre, question, bloc, durée, organisation, niveau
\newcommand{\entete}[7]{%
\begin{tcolorbox}[colback=uimm,colframe=uimm,arc=2pt,left=4pt,right=6pt,top=4pt,bottom=4pt]
\begin{minipage}[c]{18mm}\setlength{\fboxsep}{1.2pt}\fcolorbox{white}{white}{\includegraphics[height=16.5mm]{logo_uimm.jpg}}\end{minipage}\hspace{3mm}%
\begin{minipage}[c]{\dimexpr\linewidth-58mm\relax}\color{white}\raggedright
{\small Base SPC · #4}\\[1pt]{\Large\bfseries #1 · #2}\\[2pt]{\itshape #3}
\end{minipage}\hfill
\begin{minipage}[c]{33mm}\color{white}\small\raggedleft
Durée : \textbf{#5}\\ \textbf{#6}\\ Niveau #7 sur 3
\end{minipage}
\end{tcolorbox}}
\newcommand{\tagE}{{\setlength{\fboxsep}{1.5pt}\colorbox{cyan}{\color{white}\scriptsize\bfseries Ellistat}}}
\newcommand{\tagW}{{\setlength{\fboxsep}{1.5pt}\colorbox{uimm}{\color{white}\scriptsize\bfseries Page web}}}
\tcbset{pred/.style={colback=orange!6,colframe=orange!80!black,boxrule=0.6pt,arc=2pt,left=6pt,right=6pt,top=3pt,bottom=1pt,
                    title={\bfseries Je prédis, avant de commencer},coltitle=white,colbacktitle=orange!80!black,fonttitle=\small},
        retenir/.style={colback=fond,colframe=uimm,boxrule=0.6pt,arc=2pt,left=6pt,right=6pt,top=3pt,bottom=3pt,
                    title={\bfseries À retenir},coltitle=white,colbacktitle=uimm,fonttitle=\small},
        plus/.style={colback=violet!6,colframe=violet,boxrule=0.6pt,arc=2pt,left=6pt,right=6pt,top=3pt,bottom=3pt,
                    title={\bfseries Pour aller plus loin},coltitle=white,colbacktitle=violet,fonttitle=\small},
        act/.style={colback=uimm!10,colframe=uimm,boxrule=0.5pt,arc=2pt,left=5pt,right=5pt,top=2pt,bottom=2pt,before skip=7pt,after skip=3pt},
        verif/.style={colback=orange!6,colframe=orange!80!black,boxrule=0.6pt,arc=2pt,left=6pt,right=6pt,top=3pt,bottom=3pt,
                    title={\bfseries À vérifier sur poste avant la première séance},coltitle=white,colbacktitle=orange!80!black,fonttitle=\small}}
\newcommand{\fichier}[1]{\texttt{\small #1}}
\newcommand{\menu}[1]{\textsf{\textbf{#1}}}
\newcommand{\champ}[1]{\rule{#1}{0.3pt}}

\usepackage{amssymb}
\newcommand{\cb}{$\square$\,}
\newcommand{\etape}[2]{\par\vspace{9pt}{\setlength{\fboxsep}{3pt}\colorbox{uimm}{\color{white}\bfseries\large\ #1\ }}\hspace{4pt}{\color{uimm}\large\bfseries #2}\par\vspace{2pt}}
\newcommand{\motnouveau}[1]{\begin{tcolorbox}[colback=violet!6,colframe=violet,boxrule=0.6pt,arc=2pt,left=6pt,right=6pt,top=3pt,bottom=3pt,title={\bfseries Mot nouveau},coltitle=white,colbacktitle=violet,fonttitle=\small]#1\end{tcolorbox}}
\newcommand{\chiffre}[1]{\begin{tcolorbox}[colback=orange!8,colframe=orange!80!black,boxrule=1pt,arc=3pt,left=8pt,right=8pt,top=5pt,bottom=5pt]#1\end{tcolorbox}}

\usepackage{colortbl}
\pied{Bachelor Maintenance · TP 3 · Le moyen de mesure suffit-il ? · Support stagiaire}
\begin{document}
\thispagestyle{fancy}
\entete{TP 3}{Le pied à coulisse suffit-il pour réceptionner ces arbres ?}{Capabilité du moyen de mesure (Cg, Cgk, R\&R) et décision de contrôle}{Moyens de mesure · Ellistat}{3 h 30}{équipes de 3 et 4}{3}

\begin{tcolorbox}[boite]
\textbf{La situation.} Un sous-traitant livre des arbres de diamètre $\varnothing g = 48 \pm 0{,}2$ mm (dessin GUNT PT105.10.00.011). Au poste de réception, on les contrôle au pied à coulisse au 1/50 : c'est rapide. Le responsable qualité hésite à imposer le micromètre au 0,01, plus lent. Il vous demande :
\begin{enumerate}[label=\textcolor{cyan}{\textbf{\arabic*.}},leftmargin=15pt]
\item Le pied à coulisse est-il capable de contrôler cette cote ? Et le micromètre ?
\item Si le pied à coulisse ne convient pas, à partir de quelle tolérance pourrait-il servir ?
\item Quelle consigne de contrôle faut-il écrire pour le poste de réception ?
\end{enumerate}
\textbf{Ce que vous remettez.} Une note de décision d'une page (étape 7).\par
\textbf{Il vous faut.} Le coffret de cales · les 10 arbres de la mallette PT105 · un micromètre 25-50 mm et un pied à coulisse au 1/50 par équipe · Ellistat · le fichier \fichier{Releve\_RR\_arbres\_vierge.xlsx} (ordre de passage tiré au sort). Pour la cote q = 19 mm, un micromètre couvrant 0-25 mm si vous en avez un.\par
\textbf{Attention.} Les mesures sont les vôtres. Les fichiers marqués « simulé » ne servent qu'à répéter à blanc ou à remplacer une série manquante.
\end{tcolorbox}

\begin{tcolorbox}[pred]
Sur ce même arbre, le pied à coulisse (lecture au 0,02) et le micromètre (lecture au 0,01) donnent-ils des résultats répétables à mieux que 0,01 mm ? Le pied à coulisse sera-t-il jugé capable pour une tolérance de $\pm 0{,}2$ mm ? Votre réponse, avant toute mesure :\reponse[2]
\end{tcolorbox}

\titre{À faire}

\etape{1}{Quelle cote pour quel instrument ? Sans rien mesurer}
Le livret Bosch demande qu'une résolution reste au plus égale à 5 \% de la tolérance $T$ de la cote. Calculez pour chaque cote du dessin, avec une résolution de 0,02 mm (pied à coulisse) puis de 0,01 mm (micromètre) :
\begin{center}
\begin{tabularx}{0.95\linewidth}{|l|c|c|c|c|c|}
\hline
\rowcolor{fond}\textbf{Cote (dessin GUNT)} & \textbf{Écarts} & \textbf{$T$ (mm)} & \textbf{0,02 / $T$} & \textbf{0,01 / $T$} & \textbf{Instrument ?}\\\hline
$\varnothing g = 48$ (arbre) & $\pm 0{,}2$ & & & & \\\hline
$o = 34$ (plaque) & $\pm 0{,}1$ & & & & \\\hline
$n = 28$ (plaque) & $+0{,}1\ /\ -0{,}3$ & & & & \\\hline
$d = 75$ (plaque) & $+0{,}3\ /\ 0$ & & & & \\\hline
$q = 19$ (profondeur, plaque) & $0\ /\ -0{,}3$ & & & & \\\hline
$r = 16$ (profondeur, plaque) & $0\ /\ -0{,}3$ & & & & \\\hline
\end{tabularx}
\end{center}

\etape{2}{Le Cg sur une cale de 48 mm} \tagE
Composez avec le coffret une cale de $48{,}000$ mm. Notez la composition. Tenez les cales par les bords (la chaleur de la main fait varier la longueur).
\begin{enumerate}[label=\textcolor{cyan}{\textbf{\arabic*.}},leftmargin=15pt,itemsep=1.5pt]
\item Mesurez la cale \textbf{25 fois} au micromètre, puis \textbf{25 fois} au pied à coulisse. À chaque mesure, retirez l'instrument et reposez-le. Reportez les valeurs sur la fiche de relevé.
\item \menu{Data Analysis} > \menu{Gage R\&R} > \menu{Justesse}. Valeur vraie : 48{,}000 ; tolérance min : 47{,}8 ; tolérance max : 48{,}2. Une série par instrument.
\item Relevez : moyenne, écart type $s$, biais, Cg, Cgk.
\end{enumerate}
\motnouveau{$C_g = \dfrac{0{,}2\,T}{6\,s}$ et $C_{gk} = \dfrac{0{,}1\,T - |\text{biais}|}{3\,s}$. Le moyen de mesure est capable si $C_g \ge 1{,}33$ et $C_{gk} \ge 1{,}33$.}

\etape{3}{La même question sur 19 mm}
Composez $19{,}000$ mm. Même protocole, au pied à coulisse (et au micromètre si votre étendue de mesure le permet). Tolérance de la cote q : de 18,7 à 19,0 mm, soit $T = 0{,}3$ mm. Si les 25 valeurs sont toutes identiques, $s = 0$ : ce n'est pas un instrument parfait, c'est une résolution trop grosse.

\etape{4}{Le R\&R sur les arbres : 60 mesures par instrument}
Trois opérateurs (A, B, C) mesurent chacun les 10 arbres, 2 fois, \textbf{dans l'ordre tiré au sort} du fichier de relevé. Un quatrième stagiaire donne les arbres, note les valeurs et chronomètre. On fait une série complète au pied à coulisse, puis une série complète au micromètre, avec les mêmes opérateurs. Les autres équipes travaillent sur les étapes 2 et 3 pendant ce temps.
\begin{enumerate}[label=\textcolor{cyan}{\textbf{\arabic*.}},leftmargin=15pt,itemsep=1.5pt]
\item Notez la durée de chaque mesure (secondes). Ne commentez pas les valeurs avant la fin de la série.
\item Mettez les mesures au format du fichier \fichier{RR\_arbres\_*\_simule.xlsx} : une ligne par arbre, colonnes A\_r1, A\_r2, B\_r1, B\_r2, C\_r1, C\_r2.
\end{enumerate}

\etape{5}{Ellistat : le R\&R des deux instruments} \tagE
\menu{Data Analysis} > \menu{Gage R\&R} > \menu{Mesures} > créer à partir de la grille : 10 pièces, 3 opérateurs, 2 répétitions. Ressaisissez les tolérances (47,8 et 48,2) dans « Voir Plus » si elles ne sont pas conservées. Méthode ANOVA. Laissez « Production standard deviation » vide. Relevez, pour chaque instrument : répétabilité EV, reproductibilité AV, interaction, GRR, \%GRR rapporté à $T$, Cpc, nombre de catégories distinctes.

\etape{6}{Quelle tolérance minimale ?}
Pour le R\&R, la tolérance minimale contrôlable est $T_{\min} = 60 \times \text{GRR}$ (le \%GRR vaut alors 10 \%). Calculez-la pour chaque instrument et comparez à $T = 0{,}4$ mm. Pour le Cg, $T_{\min} = 39{,}9\,s + 10\,|\text{biais}|$.

\etape{7}{La note de décision}
Une page à l'attention du responsable qualité : réponse aux trois questions, chiffres à l'appui (Cg, Cgk, \%GRR, durée moyenne par mesure), consigne de contrôle proposée pour le poste de réception (instrument, nombre de mesures, tenue de l'arbre, cliquet, ordre), limites de l'étude.

\titre{J'observe, je décide}
\begin{enumerate}[label=\textcolor{cyan}{\textbf{Q\arabic*.}},leftmargin=22pt]
\item Étape 1 : quelles cotes dépassent 5 \% avec le pied à coulisse ? avec le micromètre ? Quelle cote est juste à la limite ?\reponse[2]
\item Cale 48 mm, micromètre : $s$ = \champ{14mm} \ biais = \champ{14mm} \ Cg = \champ{12mm} \ Cgk = \champ{12mm}\par\vspace{5pt}
\item Cale 48 mm, pied à coulisse : $s$ = \champ{14mm} \ biais = \champ{14mm} \ Cg = \champ{12mm} \ Cgk = \champ{12mm}\par\vspace{5pt}
\item Combien de valeurs différentes le pied à coulisse a-t-il données ? Qu'en concluez-vous sur la résolution ?\reponse[2]
\item Cale 19 mm : le pied à coulisse est-il capable pour la cote q ? Justifiez par $T$ et par le Cg.\reponse[2]
\item R\&R micromètre : EV = \champ{14mm} AV = \champ{14mm} GRR = \champ{14mm} \%GRR = \champ{12mm} Cpc = \champ{10mm} ndc = \champ{10mm}\par\vspace{5pt}
\item R\&R pied à coulisse : EV = \champ{14mm} AV = \champ{14mm} GRR = \champ{14mm} \%GRR = \champ{12mm} Cpc = \champ{10mm} ndc = \champ{10mm}\par\vspace{5pt}
\item Dans quel instrument l'écart entre opérateurs (AV) pèse-t-il le plus ? Que faut-il écrire dans la consigne pour le réduire ?\reponse[2]
\item Le livret Bosch juge le \%GRR (capable sous 10 \%, sous conditions de 10 à 30 \%, non capable au-dessus). Ellistat juge par le Cpc (rouge sous 3, orange de 3 à 4, vert au-dessus de 4). Les deux verdicts sont-ils les mêmes pour chaque instrument ? Qui fait foi pour votre note ?\reponse[3]
\item $T_{\min}$ du micromètre et du pied à coulisse. Pour quelle tolérance minimale pourrait-on utiliser chacun ?\reponse[2]
\item Durée moyenne d'une mesure : pied à coulisse \champ{14mm} s \ micromètre \champ{14mm} s. Pour 200 arbres par livraison, combien de minutes de contrôle en plus avec le micromètre ?\reponse[2]
\end{enumerate}

\begin{tcolorbox}[retenir]
Un instrument est capable \emph{pour une tolérance donnée} : le même pied à coulisse peut convenir à $\pm 1$ mm et pas à $\pm 0{,}2$ mm. Le Cg regarde un opérateur et un étalon ; le R\&R regarde des pièces réelles et plusieurs opérateurs, et l'écart entre opérateurs peut dominer. Avant de juger les pièces, on valide la mesure.
\end{tcolorbox}

\newpage
\entete{TP 3}{Fiche de relevé}{À remplir pendant la séance, puis à reporter dans Ellistat}{Moyens de mesure · Ellistat}{30 min}{à chaque poste}{3}
\vspace{4pt}
\titre{Cales : 25 mesures par instrument}
Cale : \champ{18mm} mm \quad Composition : \champ{80mm}\par\vspace{4pt}Instrument : \champ{50mm} \quad Résolution lue : \champ{14mm} mm\par\vspace{4pt}
\renewcommand{\arraystretch}{2.1}
\begin{tabularx}{\linewidth}{|*{5}{>{\centering\arraybackslash}X|}}
\hline
 1 & 2 & 3 & 4 & 5\\\hline
 6 & 7 & 8 & 9 & 10\\\hline
 11 & 12 & 13 & 14 & 15\\\hline
 16 & 17 & 18 & 19 & 20\\\hline
 21 & 22 & 23 & 24 & 25\\\hline
\end{tabularx}
\renewcommand{\arraystretch}{1.25}\vspace{4pt}
\titre{R\&R sur les arbres : instrument \champ{40mm}}
Opérateurs : A \champ{30mm} \ B \champ{30mm} \ C \champ{30mm}\par\vspace{4pt}
\begin{tabularx}{\linewidth}{|>{\centering\arraybackslash}p{14mm}|*{6}{>{\centering\arraybackslash}X|}}
\hline
\rowcolor{fond}\textbf{Arbre} & \textbf{A\_r1} & \textbf{A\_r2} & \textbf{B\_r1} & \textbf{B\_r2} & \textbf{C\_r1} & \textbf{C\_r2}\\\hline
1 & & & & & & \\\hline
2 & & & & & & \\\hline
3 & & & & & & \\\hline
4 & & & & & & \\\hline
5 & & & & & & \\\hline
6 & & & & & & \\\hline
7 & & & & & & \\\hline
8 & & & & & & \\\hline
9 & & & & & & \\\hline
10 & & & & & & \\\hline
\rowcolor{fond}\textbf{Durée (s)} & & & & & & \\\hline
\end{tabularx}
\vspace{6pt}
\begin{tcolorbox}[retenir]
Mesurez dans l'ordre tiré au sort du fichier de relevé (une série par opérateur et par répétition), pas arbre après arbre. Chaque opérateur mesure sans connaître les valeurs des autres.
\end{tcolorbox}
\end{document}
